\begin{textsample}{POS Dim 1 – vem\_ed\_18 – Score 4.00 – vem\_ed\_18/t082.txt}  \label{ex:f1_pos_116}
ALIMENTANDO COLABORAÇÕES Rosi León, diretora de Intercâmbios Virtuais e Aprendizagem Global na DePaul University (EUA), fez uma apresentação intitulada “Nurturing international faculty partnerships for COIL project implementation” (em tradução livre, “Nutrindo parcerias internacionais para implementação de projetos COIL – Collaborative Online International Learning, ou Aprendizagem Colaborativa Internacional via Online”).
Na DePaul University, os Intercâmbios Virtuais \textbf{são} chamados Global Learning Experience (GLE).
Rosi relatou o ciclo desses projetos, \textbf{que} \textbf{começam} com a divulgação das iniciativas (marketing) e o recrutamento dos professores interessados, prosseguem para a formação de parcerias, avançam para treinamentos até culminarem em seu lançamento, \textbf{desenvolvimento} e oferta em novas edições.
Ela trouxe à audiência dicas sobre como encontrar potenciais parceiros em instituições internacionais, \textbf{começando} pela rede pessoal de contatos, passando pelo coordenador COIL da instituição de ensino e incluindo redes como COIL Connect e congressos como Red LatAM COIL e IVEC.
Também \textbf{compartilhou} exemplos de aprofundamento de parcerias em projetos COIL, incluindo a visita presencial de um professor da universidade, Robert Steel, \textbf{que} irá conhecer a Fatec Tatuí na segunda quinzena de agosto de 2023.
Alejandro Molina, analista de \textbf{desenvolvimento} de carreira e mobilidade da DUOC UC Chile, abordou o tema “Espacios verdaderamente internacionales: Proyectos COIL en Portuñol”.
As experiências COIL na DUOC UC \textbf{começaram} em 2021, como internacionalização de emergência no \textbf{contexto} da pandemia de Covid-19.
Em dois anos, mais de 1000 alunos da instituição chilena \textbf{foram} atendidos em 65 Intercâmbios Virtuais, dos quais 9 com o Brasil (Centro Paula Souza, Unesp e UFSC).
As áreas de \textbf{conhecimento} abordadas nos projetos envolvem turismo, recursos naturais e gestão de negócios.
Alejandro \textbf{compartilhou} o depoimento de dois professores da DUOC UC sobre as experiências de aprendizagem e percepção sobre semelhanças e \textbf{diferenças} entre português e espanhol. “Os \textbf{idiomas} estão mais próximos do \textbf{que} percebemos, e também \textbf{é} preciso \textbf{ser} flexível para poder comunicar”, frisou Alejandro, \textbf{que} ressaltou benefícios dos Intercâmbios Virtuais entre Brasil e Chile, como “\textbf{criar} uma \textbf{ponte} entre essas culturas distintas, ainda \textbf{que} próximas”.
Recomendou aos professores interessados em projetos COIL: \textbf{ser} flexível, motivar os alunos, \textbf{construir} laços fortes com o parceiro, \textbf{aprender} sobre o país parceiro, discutir \textbf{diferenças} culturais \textbf{que} possam causar atritos entre participantes dos dois países durante o projeto e elaborar um \textbf{pequeno} glossário de palavras parecidas na escrita ou na pronúncia, mas com significados diferentes nos dois \textbf{idiomas}.

% matched lemmas: aprender, começar, compartilhar, conhecimento, construir, contexto, criar, desenvolvimento, diferença, idioma, pequeno, ponte, que, ser
\end{textsample}
